\documentclass[11pt]{article}
\usepackage[margin=1in]{geometry}
\usepackage{amsmath,amssymb,amsthm}
\usepackage{enumitem}

\newtheorem{theorem}{Theorem}
\newtheorem{lemma}{Lemma}
\newtheorem{corollary}{Corollary}
\theoremstyle{definition}
\newtheorem{remark}{Remark}

\DeclareMathOperator{\Imag}{Im}
\newcommand{\vtwo}{v_2}
\newcommand{\ff}[2]{(#1)_{#2}}

\title{The two-row $d=4$ law for all $b\equiv 0\pmod 4$:\\ a $2$-adic valuation identity}
\author{Clio}
\date{2026-06-08}

\begin{document}
\maketitle

\begin{abstract}
We close the case $b\equiv 0\pmod 4$ of the two-row $d=4$ fiber-vanishing law
$G_{(2m-b,b)}(i)=0\iff (2,2)$. Together with the previously proved class
$b\equiv 1\pmod 4$, this establishes the law unconditionally for the full residue
classes $b\equiv 0,1\pmod 4$ --- half of all $b$. The proof is the exact $2$-adic
identity $\vtwo(I_b(m))=\vtwo(\ff{m}{R+1})-\vtwo(R!)$, valid for every integer $m$.
In contrast with the $b\equiv 1$ case --- where a single term of the binomial
expansion strictly dominates --- here, for odd $m$, exactly three terms attain the
minimal valuation; the law survives because \emph{three odd numbers sum to an odd
number}. The parity of the number of valuation-minimal terms is the operative
invariant.
\end{abstract}

\section{Setup (proved in prior cycles; not re-derived)}

Fix $\lambda=(2m-b,b)$ with $0\le b\le m$, put $s=1+i$, $P(u)=(1+su+u^2)^m$, and
\[
   I_b(m):=\Imag G_{(2m-b,b)}=\Imag\,[u^b]\bigl((1-u)P(u)\bigr)\in\mathbb Z .
\]
Earlier work establishes the following, which we take as given.

\begin{itemize}[leftmargin=1.4em]
\item \textbf{Exact expansion (E).} Writing $1+su+u^2=(1+u^2)+su$ and extracting
imaginary parts,
\[
   I_b(m)=\sum_{j=1}^{b}\tau_j(m),\qquad
   \tau_j(m)=\binom{m}{j}\,\Imag\bigl((1+i)^j\bigr)\,\bigl(T(b-j)-T(b-1-j)\bigr),
\]
where $T(\ell)=\binom{m-j}{\ell/2}$ for $\ell$ even and $T(\ell)=0$ for $\ell$ odd.
Of the two extractions exactly one is nonzero (their indices $b-j,\,b-1-j$ are
consecutive), so $T(b-j)-T(b-1-j)=\pm\binom{m-j}{\beta_j}$ with
$\beta_j:=\lfloor (b-j)/2\rfloor$.

\item \textbf{The complex factor.} $\Imag\bigl((1+i)^j\bigr)=2^{j/2}\sin(j\pi/4)\in\mathbb Z$, with
\[
   \vtwo\Bigl(\Imag\bigl((1+i)^j\bigr)\Bigr)=\Big\lfloor \tfrac j2\Big\rfloor\quad(j\not\equiv 0\bmod 4),
   \qquad \Imag\bigl((1+i)^j\bigr)=0\quad(j\equiv 0\bmod 4).
\]
In particular $\Imag((1+i)^j)$ is odd iff $j=1$.

\item \textbf{Forced roots and the reduction.} $I_b(m)$ is an integer-valued
polynomial in $m$ whose only integer roots are the forced ones
$m\in\{0,1,\dots,R\}$, $R:=\lfloor (b-1)/2\rfloor$. Writing
$I_b(m)=\ff{m}{R+1}\,Q_b(m)$ with $\ff{m}{R+1}=\prod_{r=0}^{R}(m-r)$, the two-row
law is exactly
\[
   (\star)\qquad Q_b(m)\neq 0\ \text{ for every integer } m\ge b .
\]
Since $m\ge b>R$ makes $\ff{m}{R+1}\neq 0$, $(\star)$ is equivalent to
$I_b(m)\neq 0$ for $m\ge b$.
\end{itemize}

Throughout this note $b\equiv 0\pmod 4$, say $b=4t$. Then
\[
   R=\frac{b-2}2=2t-1\ \text{ is \emph{odd}},\qquad R+1=\frac b2\ \text{ is even}.
\]
The oddness of $R$ is the structural fact that drives everything below.

\medskip
\noindent\textbf{Main result.}
\begin{theorem}\label{thm:main}
Let $b\equiv 0\pmod 4$ and $R=(b-2)/2$. For every integer $m$,
\[
   (\mathrm V)\qquad \vtwo\bigl(I_b(m)\bigr)=\vtwo\bigl(\ff{m}{R+1}\bigr)-\vtwo(R!).
\]
Consequently $\vtwo(Q_b(m))=-\vtwo(R!)$ is constant, so $Q_b(m)\neq 0$ for every
integer $m$; in particular $(\star)$ holds and the two-row $d=4$ law holds for all
$b\equiv 0\pmod 4$.
\end{theorem}

The same identity $(\mathrm V)$ was proved last cycle for $b\equiv 1\pmod 4$, so
\textbf{the two-row $d=4$ law now holds for all $b\equiv 0,1\pmod 4$.}

\section{The $j=1$ term and the abbreviation $a$}

For $b$ even, $b-1$ is odd and $b-2$ is even, so in $\tau_1$ the surviving
extraction is $T(b-2)=\binom{m-1}{R}$ (here $\beta_1=\lfloor(b-1)/2\rfloor=R$),
giving
\[
   \tau_1(m)=-\,m\binom{m-1}{R}=-(R+1)\binom{m}{R+1}=-\frac{\ff{m}{R+1}}{R!}.
\]
Hence, abbreviating
\[
   a:=\vtwo\bigl(\tau_1(m)\bigr)=\vtwo\bigl(\ff{m}{R+1}\bigr)-\vtwo(R!),
\]
the right-hand side of $(\mathrm V)$ is exactly $a$. The theorem is the statement
$\vtwo(I_b(m))=a$.

\section{Term valuations: an exact decomposition}\label{sec:decomp}

Fix $j$ with $\tau_j\neq 0$, i.e. $j\not\equiv 0\pmod 4$, and put
$h:=\lfloor j/2\rfloor$. We compute $\vtwo(\tau_j)-a$ exactly.

Since $b$ is even, $\beta_j=\lfloor (b-j)/2\rfloor$ and one checks
$s_j:=j+\beta_j=R+1+h$ in both parities of $j$ (for $j=2h$: $\beta_j=R+1-h$; for
$j=2h+1$: $\beta_j=R-h$). The Vandermonde identity
$\binom{m}{j}\binom{m-j}{\beta_j}=\binom{m}{s_j}\binom{s_j}{j}$ gives
\[
   \vtwo(\tau_j)=\vtwo\!\binom{m}{s_j}+\vtwo\!\binom{s_j}{j}+\Big\lfloor\tfrac j2\Big\rfloor .
\]
Telescoping $\dfrac{\binom{m}{s_j}}{\binom{m}{R+1}}
=\prod_{t=1}^{h}\dfrac{m-R-t}{R+1+t}$ (note $s_j-(R+1)=h$) and subtracting
$a=\vtwo(R+1)+\vtwo\binom{m}{R+1}$ yields, after cancelling the common
$\vtwo\binom{m}{R+1}$,
\[
   \boxed{\ \vtwo(\tau_j)-a \;=\; D_0(j)+S_j(m),\qquad
   S_j(m):=\sum_{t=1}^{h}\vtwo(m-R-t),\ }
\]
where the $m$-independent offset is
\[
   D_0(j)=h+\vtwo\!\binom{s_j}{j}-\sum_{t=1}^{h}\vtwo(R+1+t)-\vtwo(R+1)
        =h+\vtwo\!\binom{s_j}{j}-\bigl(\vtwo((R+1+h)!)-\vtwo(R!)\bigr).
\]
Using $s_j=R+1+h$, so $\binom{s_j}{j}=\dfrac{(R+1+h)!}{j!\,\beta_j!}$, the factor
$(R+1+h)!$ cancels and $D_0(j)=h+\vtwo(R!)-\vtwo(j!)-\vtwo(\beta_j!)$. Evaluating
the two parities with $\vtwo((2h)!)=\vtwo((2h+1)!)=h+\vtwo(h!)$:

\begin{lemma}\label{lem:D0}
For $\tau_j\neq 0$ and $h=\lfloor j/2\rfloor$,
\[
   D_0(j)=
   \begin{cases}
     \vtwo\!\dbinom{R}{h}, & j\ \text{odd},\\[1.2ex]
     \vtwo\!\dbinom{R}{h-1}-\vtwo(h), & j\ \text{even}.
   \end{cases}
\]
\end{lemma}

\begin{proof}
\emph{$j=2h+1$ odd:} $\beta_j=R-h$, so
$D_0=h+\vtwo(R!)-\bigl(h+\vtwo(h!)\bigr)-\vtwo((R-h)!)
=\vtwo(R!)-\vtwo(h!)-\vtwo((R-h)!)=\vtwo\binom{R}{h}$.

\emph{$j=2h$ even:} $\beta_j=R+1-h$, so
$D_0=h+\vtwo(R!)-\bigl(h+\vtwo(h!)\bigr)-\vtwo((R+1-h)!)
=\vtwo(R!)-\vtwo(h!)-\vtwo((R+1-h)!)$.
Since $(R+1-h)!=(R-h)!\,(R+1-h)$ we get
$D_0=\vtwo\binom{R}{h}-\vtwo(R+1-h)$. Finally the elementary identity
$h\binom{R}{h}=(R+1-h)\binom{R}{h-1}$ --- both sides equal
$\tfrac{R!}{(h-1)!\,(R-h)!}$ --- gives, on taking $\vtwo$,
$\vtwo\binom{R}{h}-\vtwo(R+1-h)=\vtwo\binom{R}{h-1}-\vtwo(h)$.
\end{proof}

\section{No term dips below $\tau_1$}

\begin{lemma}[(P1)]\label{lem:P1}
For every $j$ with $\tau_j\neq0$ and every integer $m$, $\vtwo(\tau_j)\ge a$, i.e.
$D_0(j)+S_j(m)\ge 0$.
\end{lemma}

\begin{proof}
For $j$ odd, $D_0(j)=\vtwo\binom{R}{h}\ge0$ and $S_j\ge0$, done.

For $j$ even, the $h$ integers $m-R-1,m-R-2,\dots,m-R-h$ are consecutive, so their
product is divisible by $h!$; hence
\[
   S_j(m)=\vtwo\!\Bigl(\textstyle\prod_{t=1}^{h}(m-R-t)\Bigr)\ \ge\ \vtwo(h!)\ \ge\ \vtwo(h).
\]
By Lemma~\ref{lem:D0}, $D_0(j)=\vtwo\binom{R}{h-1}-\vtwo(h)\ge -\vtwo(h)$. Adding,
$D_0(j)+S_j\ge -\vtwo(h)+\vtwo(h)=0$.
\end{proof}

\section{Exactly which terms tie}

\begin{lemma}[tie classification]\label{lem:ties}
Among $j\ge 2$ with $\tau_j\neq 0$, the equality $\vtwo(\tau_j)=a$ holds for some
$m$ only if $j\in\{2,3\}$. Moreover, for $R$ odd:
\[
   \vtwo(\tau_2)=\vtwo(\tau_3)=a \iff m\ \text{odd},\qquad
   \vtwo(\tau_2),\vtwo(\tau_3)>a \iff m\ \text{even},
\]
and for every $j\ge 4$ with $\tau_j\neq 0$ one has $\vtwo(\tau_j)>a$ for all $m$.
\end{lemma}

\begin{proof}
By Lemma~\ref{lem:P1} a tie means $D_0(j)+S_j(m)=0$. We bound $S_j$ from below.

\emph{$j$ even, $h\ge 3$:} $S_j\ge \vtwo(h!)=\vtwo(h)+\vtwo((h-1)!)\ge\vtwo(h)+1$,
because $(h-1)!$ is even for $h\ge 3$. With $D_0(j)\ge-\vtwo(h)$ this gives
$D_0(j)+S_j\ge 1>0$: strict.

\emph{$j$ even, $h=2$:} then $j=4\equiv 0\pmod 4$, so $\tau_4=0$ --- excluded.

\emph{$j$ odd, $h\ge 2$ (i.e. $j\ge 5$):} the list $m-R-1,\dots,m-R-h$ contains two
consecutive integers, one even, so $S_j\ge 1$; with $D_0(j)\ge0$, $D_0(j)+S_j\ge1>0$:
strict.

Thus the only $j\ge2$ that can tie are $j=2$ ($h=1$) and $j=3$ ($h=1$). For both,
$S_j=\vtwo(m-R-1)$, and by Lemma~\ref{lem:D0}, using $R$ odd,
\[
   D_0(2)=\vtwo\!\binom{R}{0}-\vtwo(1)=0,\qquad
   D_0(3)=\vtwo\!\binom{R}{1}=\vtwo(R)=0 .
\]
Hence $\vtwo(\tau_2)-a=\vtwo(\tau_3)-a=\vtwo(m-R-1)$. Since $R$ is odd, $m-R-1$ is
odd iff $m$ is odd; therefore both differences vanish iff $m$ is odd, and are $\ge1$
iff $m$ is even.
\end{proof}

\section{Proof of the Theorem}

By Lemma~\ref{lem:P1}, $\vtwo(\tau_j)\ge a$ for all $j$, so
$N(m):=I_b(m)/2^{a}=\sum_{j}\tau_j(m)/2^{a}$ is an integer, and modulo $2$ only the
\emph{tying} terms (those with $\vtwo(\tau_j)=a$) survive:
\[
   N(m)\equiv \sum_{j:\,\vtwo(\tau_j)=a}\frac{\tau_j(m)}{2^{a}}\pmod 2 .
\]
Each summand $\tau_j/2^{a}$ is odd (its $2$-adic valuation is exactly $0$). By
Lemma~\ref{lem:ties}:

\begin{itemize}[leftmargin=1.4em]
\item If $m$ is even, the only tying term is $j=1$, so $N(m)\equiv 1\pmod 2$.
\item If $m$ is odd, the tying terms are exactly $j\in\{1,2,3\}$, so $N(m)$ is a sum
of \emph{three} odd integers, hence $N(m)\equiv 1\pmod 2$.
\end{itemize}

In both cases $N(m)$ is odd, i.e. $\vtwo(I_b(m))=a=\vtwo(\ff{m}{R+1})-\vtwo(R!)$,
which is $(\mathrm V)$. For $m\ge b>R$ the right-hand side is finite, so
$I_b(m)\neq0$; equivalently $Q_b(m)\neq0$. This is $(\star)$, the two-row $d=4$ law
for $b\equiv 0\pmod4$. \qed

\begin{remark}
The mechanism differs sharply from $b\equiv1\pmod4$. There $R+1$ is odd, every
$\tau_j$ with $j\ge2$ satisfies $\vtwo(\tau_j)>a$ strictly, and $I_b/2^a$ reduces to
the single odd term $\tau_1/2^a$. For $b\equiv0\pmod4$ the factor $R+1$ is even,
which lowers the offsets so that, for odd $m$, the terms $j=2,3$ \emph{tie} $\tau_1$.
The law is rescued not by domination but by counting: the number of
valuation-minimal terms is odd ($1$ for even $m$, $3$ for odd $m$), so their sum
keeps an odd leading $2$-adic digit. The oddness of $R$ --- a direct consequence of
$4\mid b$ --- is what forces both $D_0(2)=D_0(3)=0$ and pins the tie to the parity
of $m$.
\end{remark}

\begin{remark}
Equivalently, over $\mathbb F_2$ the monic irreducible part $q_b$ of $Q_b$ satisfies
$q_b(m)\equiv (m^2+m+1)^{\lfloor(b-1)/4\rfloor}\pmod2$ for $b\equiv0,1\pmod4$
(verified $b\le40$); since $m^2+m+1$ is odd at every integer this re-expresses
``no integer root.'' Theorem~\ref{thm:main} proves the underlying non-vanishing
directly via $(\mathrm V)$, without needing the polynomial factorization shape.
\end{remark}

\section*{Verification}
The decomposition of \S\ref{sec:decomp}, the closed forms of Lemma~\ref{lem:D0},
the no-dip bound (P1), the tie classification, and the identity $(\mathrm V)$ were
checked exactly for all $b\equiv0\pmod4$ with $4\le b\le 64$, over
$m\in[b,b+20]$ together with $40$ random $m\le 10^6$ per $b$ ($24{,}480$ term-level
checks; \texttt{verify.py}: ``ALL OK'').

\end{document}
